\documentclass[a4paper]{article}
% Shared preamble for the English probability and statistics documents.
% Based on the original course preamble. Compile with pdfLaTeX.
\usepackage[top=1.5cm,bottom=2cm,left=2cm,right=2cm]{geometry}
\usepackage[T1]{fontenc}
\usepackage{lmodern}
\usepackage[utf8]{inputenc}
\usepackage[english]{babel}

\usepackage{amsmath}
\usepackage{amssymb}
\usepackage{amsthm}
\usepackage{mathrsfs}
\usepackage{mathtools}
\usepackage{etoolbox}
\usepackage{tikz}
\usetikzlibrary{shapes.multipart,arrows.meta}
\usepackage{xcolor}
\usepackage{enumerate}
\usepackage[shortlabels]{enumitem}
\usepackage{hyperref}
\usepackage{version}
\usepackage{bbold}
\usepackage{fontawesome5}
\usepackage{fancyhdr}
\lfoot{}
\usepackage{natbib}
\usepackage{graphicx}
\usepackage{pdfpages}
\usepackage{authblk}
    
\usepackage{blindtext}
\usepackage{multicol}
\usepackage{collcell}
\usepackage{wrapfig}
\usepackage{blkarray}
\usepackage{cancel}
\usepackage{appendix}
\usepackage{float}
\makeatletter
\input{tex/assets/slashbox.sty}
\makeatother
\hypersetup{unicode,colorlinks,citecolor=red,filecolor=black,
  linkcolor=red,urlcolor=cyan}
\setlength{\columnsep}{1cm}
\pagenumbering{gobble}

\def\numtext#1{%
  \ifcase#1\or one\or two\or three\or four\or five\or six\or seven\or eight\or nine\or ten\or
  eleven\or twelve\or thirteen\or fourteen\or fifteen\or sixteen\or seventeen\or eighteen\or nineteen\or twenty\or Lots\fi}

\newcommand{\Byte}[1]{%
  \def\nodecontents{}%
  \foreach[count=\l] \k in {#1} {\xappto{\nodecontents}{\noexpand\nodepart{\numtext{\l}}\k}}%
  \node[rectangle split,rectangle split parts=8,rectangle split horizontal,draw] at (2,2){\nodecontents};}

\newcommand{\D}{{\mathrm{d}}}
\newcommand{\cE}{{\mathcal{E}}}
\newcommand{\cG}{{\mathcal{G}}}
\newcommand{\cX}{{\mathcal{X}}}
\newcommand{\cY}{{\mathcal{Y}}}
\newcommand{\cB}{{\mathcal{B}}}
\newcommand{\cF}{{\mathcal{F}}}
\newcommand{\cO}{{\mathcal{O}}}
\newcommand{\cS}{{\mathcal{S}}}
\newcommand{\cT}{{\mathcal{T}}}
\newcommand{\given}{\, \big| \,}
\newcommand{\cL}{{\mathscr{L}}}
\newcommand{\cD}{{\mathcal{D}}}
\newcommand{\cA}{{\mathscr{A}}}
\newcommand{\cP}{{\mathscr{P}}}
\newcommand{\cR}{{\mathscr{R}}}
\newcommand{\A}{{\mathcal A}}
\newcommand{\one}{\mathbb{1}}
\newcommand{\mathbbm}[1]{\mathbb{#1}}

\providecommand{\C}{}
\renewcommand{\C}{\mathbb C}
\newcommand{\R}{\mathbb R}
\newcommand{\Q}{\mathbb Q}
\newcommand{\N}{\mathbb N}
\newcommand{\T}{\mathbb T}
\newcommand{\F}{\mathbb F}
\renewcommand{\P}{\mathbb P}
\newcommand{\Z}{\mathbb Z}
\newcommand{\E}{\mathbb {E}}
\newcommand{\Pbb}[1]{\P\left( #1\right)}
\newcommand{\Ebb}[1]{\E\left[ #1\right]}

\newcommand{\lb}{\left (}
\newcommand{\rb}{\right )}
\newcommand{\lbb}{\left [}
\newcommand{\rbb}{\right ]}
\newcommand{\labs}{\left |}
\newcommand{\rabs}{\right |}
\newcommand{\lbrb}[1]{\lb #1 \rb}
\newcommand{\lbbrbb}[1]{\lbb#1\rbb}
\newcommand{\labsrabs}[1]{\labs#1\rabs}
\newcommand{\lbbrb}[1]{\lbb#1\rb}
\newcommand{\lbrbb}[1]{\lb#1\rbb}
\newcommand{\langlerangle}[1]{\langle#1\rangle}
\newcommand{\lbcurly}{\{}
\newcommand{\rbcurly}{\}}
\newcommand{\curly}[1]{\lbcurly#1\rbcurly}
\newcommand{\abs}{\labsrabs}
\newcommand{\ind}[1]{1_{#1}}
\newcommand{\indep}{\perp \!\!\! \perp}
\renewcommand{\d}{\, \mathrm{d}}

\newtheorem{theorem}{Theorem}
\newtheorem{problemCnt}{Lemma}
\theoremstyle{definition}
\newtheorem{problem}[problemCnt]{Problem}
\newtheorem{problembook}[problemCnt]{\faBook \,  Problem}
\newtheorem{problempencil}[problemCnt]{\faPencil* \,  Problem}
\newtheorem{problemstar}[problemCnt]{\faStar\,  Problem}
\newtheorem{problemstarpencil}[problemCnt]{\faStar \, \faPencil* \,   Problem}
\newtheorem{probleminteresting}[problemCnt]{\faSmile[regular] \,   Problem}
\newtheorem{solution}{Solution}

\newcommand{\fwcell}[1]{\makebox[\arraystretch\normalbaselineskip]{$#1$}}
\newenvironment{enumeratei}{\begin{enumerate}[(i)]}{\end{enumerate}}
\newenvironment{enumeratea}{\begin{enumerate}[(a)]}{\end{enumerate}}
\newenvironment{enumeraten}{\begin{enumerate}[\upshape 1.]}{\end{enumerate}}

% Compatibility helpers needed by a few files in the combined archive.
\AtBeginDocument{\noindent{\footnotesize English translation prepared with ChatGPT (OpenAI, GPT-6), 26 October 2026, for translation of the Bulgarian text into English and LaTeX formatting.
The Bulgarian originals remain the reference; translation may contain errors.}\par\smallskip}
\newcommand{\MissingGraphic}[1]{%
  \begin{center}\fbox{\parbox{0.82\linewidth}{\centering Missing source image: \texttt{#1}}}\end{center}}



\setlength{\emergencystretch}{3em}
\pagenumbering{arabic}
\begin{document}
\begin{center}\Large Probability and Statistics, Computer Science\par\large Regular exam, 2 July 2023: solutions\end{center}
\noindent\textit{English transcription of the supplied handwritten workings. The source's numerical conclusions and noted ambiguities are retained.}
\section*{Problem 1}
\subsection*{1. Percentage measurement error}
\[\P(\text{error}>10\%)=\P(N(0,5^2)>10)=\P(N(0,1)>2)=1-\Phi(2)=:p\approx2.28\%.\]
The expected number of cars until the first such error is $\E[Ge(p)]=1/p$. Among 1000 cars, the expected number with error above 10\% is $\E[Bin(1000,p)]=1000p$.
For iid $X_i\sim N(0,5^2)$,
\begin{align*}
\P\left(\frac{X_1+\cdots+X_{100}}{100}>1\right)
&=\P(N(0,100\cdot5^2)>100)\\
&=\P(N(0,1)>2)=1-\Phi(2)\approx2.28\%.
\end{align*}
\subsection*{2. Car speed and recorded speed}
Let $V\sim N(40,5^2)$ be speed and $X\sim N(0,5^2)$ instrument error. Then $\P(V>50)=\P(N(0,1)>2)\approx2.28\%$.
The notes point out that the actual question asks for the probability that the instrument reports more than 50 km/h:
\[\P\bigl(V(1+X/100)>50\bigr)=\P\left(N_1(40,5^2)\,N_2\left(1,(1/20)^2\right)>50\right).\]
It is difficult to evaluate explicitly; either $\P(V>50)$ or the expression above was accepted. A Monte Carlo estimate using $10^6$ simulations is about 3.46\%.

For the mean absolute speed error represented by signed products, the notes seek $\sigma$ from
\[\P\left(\frac{X_1V_1+\cdots+X_{100}V_{100}}{10000}>1\right)\le0.1\%.\]
Independence gives
\[\E(X_1V_1)=\E X_1\E V_1=0,\qquad
D(X_1V_1)=\E X_1^2\E V_1^2-(\E X_1)^2(\E V_1)^2=\sigma^2(40^2+25).\]
The CLT approximation is
\[\P\left(N(0,1)>\frac{10000}{\sigma\sqrt{1625}\sqrt{100}}\right).\]
The handwritten calculation sets the threshold approximately equal to 3.72 and obtains $\sigma\approx6.67$. It comments that the exam wording was ambiguous and accepts any correct calculation with a comparable interpretation.
\subsection*{3. Absolute errors}
For $|N(0,5^2)|$, the notes use $p=2\Phi(-2)$ and repeat the geometric and binomial arguments. For the average absolute error they write
\[\P\left(\frac{|X_1|+\cdots+|X_{100}|}{100}>1\right)
\approx\P\left(N(0,1)>\frac{100-100\E|X_1|}{\sqrt{100D|X_1|}}\right).\]
The displayed moment calculation is
\[\E|X_1|=5\E|N(0,1)|=5\cdot2\int_0^\infty\frac{x}{\sqrt{2\pi}}e^{-x^2/2}\,dx=\frac5{\sqrt{2\pi}},\]
followed by $D|X_1|=25-25/(2\pi)$ and ``substitute''.
\noindent\textit{Transcription note: the last equality for the absolute normal mean drops a factor of 2 in the handwritten source. Its following variance and numerical threshold above are preserved, rather than silently replaced.}
\section*{Problem 2}
\subsection*{1. Independence of the counts}
Let $I=i_1+\cdots+i_n$ and $F_n=f(1)+\cdots+f(n)$. Conditioning on $N$ gives
\begin{align*}
&\P(X_1=i_1,\dots,X_n=i_n)\\
&=\sum_{k=I}^\infty\P(N=k)\P(X_1=i_1,\dots,X_n=i_n\mid N=k)\\
&=\sum_{k=I}^\infty\frac{e^{-\lambda}\lambda^k}{k!}
\frac{k!}{i_1!\cdots i_n!(k-I)!}\prod_{j=1}^nf(j)^{i_j}(1-F_n)^{k-I}\\
&=e^{-\lambda}\prod_{j=1}^n\frac{(\lambda f(j))^{i_j}}{i_j!}
\sum_{r=0}^\infty\frac{\lambda^r(1-F_n)^r}{r!}\\
&=\prod_{j=1}^n\frac{e^{-\lambda f(j)}(\lambda f(j))^{i_j}}{i_j!}.
\end{align*}
Thus the $X_i$ are independent and $X_i\sim Poi(\lambda f(i))$.
\subsection*{2. Expected sales until stopping}
The event $\{T\le n-1\}$ depends only on $X_1,\dots,X_{n-1}$, so it is independent of $X_n,X_{n+1},\dots$. Its complement $\{T\ge n\}$ is therefore independent of $X_n$. Consequently,
\begin{align*}
\E S&=\E\sum_{i=1}^\infty X_i\one_{\{T\ge i\}}
=\sum_{i=1}^\infty\E X_i\P(T\ge i)\\
&=\sum_{i=1}^\infty\E X_i\sum_{j=i}^\infty\P(T=j)
=\sum_{j=1}^\infty\P(T=j)\sum_{i=1}^j\E X_i\\
&=\lambda\sum_{j=1}^\infty\P(T=j)(f(1)+\cdots+f(j))
=\lambda\E[f(1)+\cdots+f(T)].
\end{align*}
\textit{Warning in the source:} starting instead from
\[\E S=\sum_{k\ge1}\P(T=k)\E[X_1+\cdots+X_k\mid T=k],\]
one cannot replace the conditional expectation by $\E X_1+\cdots+\E X_k$, since $\{T=k\}$ is an event determined by $X_1,\dots,X_k$ and is not independent of them.
\section*{Problem 3}
\subsection*{1. Fractional parts}
Let $X=\{MU\}$. For $t\in(0,1)$,
\[\P(X\le t)=\sum_{k=0}^{M-1}\P(MU\in[k,k+t])=M\frac tM=t,\]
since $MU\sim U(0,M)$. Hence $X\sim U(0,1)$, $\E X=1/2$ and $DX=1/12$.
To compute correlation,
\[Cor(X,U)=\frac{\E[XU]-1/4}{1/12}.\]
Partitioning the unit interval,
\begin{align*}
\E[XU]&=\sum_{k=0}^{M-1}\int_{k/M}^{(k+1)/M}(Mt-k)t\,dt\\
&=M\int_0^1t^2dt-\sum_{k=1}^{M-1}k\int_{k/M}^{(k+1)/M}t\,dt\\
&=\frac M3-\frac1{2M^2}\sum_{k=1}^{M-1}(2k^2+k)\\
&=\frac M3-\frac1{2M^2}\left[2\frac{(M-1)M(2M-1)}6+\frac{(M-1)M}2\right]\\
&=\frac1{12M}\bigl[4M^2-(M-1)(4M+1)\bigr]
=\frac1{12}\left(3+\frac1M\right).
\end{align*}
Thus $Cor(X,U)=1/M$. The larger $M$, the less correlated $\{MU\}$ and $U$ are.
\subsection*{2. Extremes of 100 uniforms}
Each $U_i$ has mean $1/2$. For $0<t<1$,
\[\P(\max(U_1,\dots,U_{100})\le t)=\P(U_1\le t)^{100}=t^{100}.\]
The maximum has density $100t^{99}$ and expectation $\int_0^1100t^{100}dt=100/101$. Since $1-U_i\overset d=U_i$,
\[\min(U_1,\dots,U_{100})\overset d=1-\max(U_1,\dots,U_{100}),\]
so its expectation is $1/101$. Alternatively, derive its distribution function directly.
\subsection*{3. Tangent of a uniform angle}
For $V\sim U(-\pi/2,\pi/2)$ and $Y=\tan V$,
\[\P(Y\le t)=\P(V\le\arctan t)=\int_{-\pi/2}^{\arctan t}\frac1\pi dx=\frac{\arctan t+\pi/2}{\pi}.\]
Hence $f_Y(t)=1/(\pi(1+t^2))$ for $t\in\R$. The integrand $t/(\pi(1+t^2))$ behaves as $1/(\pi t)$ at infinity and is not integrable, so $\E Y$ and $DY$ are not defined.
\subsection*{4. Cauchy transformations}
Write $Z\overset d=\tan V$ as above. Then
\[1/Z\overset d=1/\tan V=\cot V\overset d=\tan V,\]
with the instruction ``formalise why!''. Also
\[\frac{2\tan V}{1-\tan^2V}=\tan(2V)\overset d=\tan V,\qquad
\frac{3\tan V-\tan^3V}{1-3\tan^2V}=\tan(3V)\overset d=\tan V.\]
The usual density-transformation method is another option.
\section*{Bonus}
\begin{enumerate}
\item See the law of total variance.
\item See Wikipedia; by symmetry $\E X=1/2$, and $DX=1/8$.
\item The derivative is zero almost everywhere, so $X$ has no density. The Cantor function is an example of a continuous increasing function that is not absolutely continuous.
\end{enumerate}
\end{document}
